%% =====================================================================
%% (1) main.tex preamble: add next to \MsevenEight, \MsevenNine and \MsevenTen.
%% SymPort's total QTM over all 184 Santa puzzles with N = 4..10:
%%   6,334 (N = 4..7, already measured) + 3*\MsevenEight + 3*\MsevenNine + 4*\MsevenTen
%% KMCoders' published total over the same 184 puzzles is 12,324.
\newcommand{\MsevenSantaTotal}{??,???}

%% =====================================================================
%% (2) main.tex: replace everything between \begin{abstract} and \end{abstract}
\begin{abstract}
Short solutions of large Rubik's puzzles are known mostly from small cubes and from programming competitions, and the methods
are rarely compared on common ground. We build one open engine for the \Cu{N} Rubik's Cube and the \Sq{N} Rubik's Square and
run eight solvers on the same instances. Five reproduce known methods, among them the KMCoders pipeline behind the strongest
public large-cube solutions of Kaggle's Santa 2023 competition, and three are new. Our main method, \Mseven, runs a deterministic
solver on symmetric copies of a scramble (up to 48 for the cube, 8 for the Square) and keeps the shortest answer, which maps back
to a solution of the original state with the same length. Wrapped around KMCoders' own program with an optimal central
3$\times$3$\times$3 stage, it is shorter on average than their published best-of-many solutions at every Santa 2023 cube size
from $N=4$ to 10, with \MsevenSantaTotal{} against 12,324 quarter turns over all \SantaTotal{} puzzles. \FiveC schedules
Demaine et al.'s bulk commutators by net gain rather than unanimous need. On random cubes it is \FiveCGainTen\% shorter than
Demaine's algorithm at $N=10$, a lead that grows with $N$. \FiveA, a cluster-token transformer trained on small Squares only,
solves Squares of sizes it never saw. Finally, a counting lower bound, calibrated on Squares whose exact God's numbers we
compute (\GodSquareThree, \GodSquareFour{} and \GodSquareFive{} for $N=3$, 4 and 5), puts the published Santa solutions at
\PubOverBoundFour{}$\times$ the bound at $N=4$ and \PubOverBoundTen{}$\times$ at $N=10$, a gap that grows like $\log N$. We
release the engine, the benchmarks, every solver and every result file.
\end{abstract}

%% =====================================================================
%% (3) sections/intro.tex: replace the whole file (the separate Contributions section is folded in here)
%% =====================================================================
\section{Introduction}
\label{sec:intro}

The \Cu{N} Rubik's Cube is a standard testbed for planning on graphs that are far too large to
enumerate. Its worst-case solution length, God's number, is 20 face turns for $N=3$~\cite{rokicki2014diameter}
(26 quarter turns~\cite{rokicki2014qtm}) and grows as $\Theta(N^2/\log N)$~\cite{demaine2011algorithms};
finding a shortest solution of a given state is NP-complete~\cite{demaine2018npcomplete}. Learned solvers
have recently become strong on small cubes. DeepCubeA~\cite{agostinelli2019deepcubea} and
EfficientCube~\cite{takano2023selfsupervision} solve the \Cu{3}, and CayleyPy~\cite{chervov2025cayleypy}
reaches near-optimal \Cu{3} solutions and the first learned \Cu{4} and \Cu{5} solvers. All of them train one
network per size, and none reports a cube beyond $N=5$. At the other end of the scale, the strongest public
large-cube solutions come from a programming competition, Kaggle's Santa 2023~\cite{kaggle2023santa}, where
the team KMCoders solved, among others, a \Cu{33} in 4,719 quarter turns~\cite{kmcoders2024santa}.

We ask what it takes to cover the whole range of sizes, for both the cube and its flat relative, the \Sq{N}
Rubik's Square of Demaine et al.~\cite{demaine2011algorithms}; which approach gives the shortest solutions at
each size; whether any method can beat the published competition solutions; and how far all of them are
from optimal. To answer on common ground, we run eight solvers and three external baselines on one open engine,
on the same seeded random Squares and the same Santa 2023 cube puzzles; every number in this paper is generated
by script from the result files.

\paragraph{Contributions.}
\begin{itemize}
\item \textbf{\Mseven, a symmetry portfolio} (\S\ref{sec:m7}), our main method. Around KMCoders' program it is the best
  method at every Santa 2023 cube size from $N=4$ to 10, ahead of their published best-of-many solutions on average
  (Table~\ref{tab:maincube}); around our strongest Square solvers it is the best method at \MsevenRedSquare{} of the
  \SizesSquare{} Square sizes (Table~\ref{tab:mainsq}).
\item \textbf{\FiveC, net-gain scheduling of bulk commutators} (\S\ref{sec:5c}): the strongest single solver on large
  Squares, and \FiveCGainTwenty\% shorter than Demaine's algorithm on the random \Cu{20}.
\item \textbf{\FiveA, one network for every size} (\S\ref{sec:5a}): a cluster-token transformer that solves Square
  sizes it never saw in training.
\item \textbf{How far from optimal} (\S\ref{sec:bound}): a counting lower bound for every \Cu{N}, calibrated on the
  first exact God's numbers we know of for the \Sq{4} and \Sq{5}, and measured against every method.
\item \textbf{An open engine and benchmark} (\S\ref{sec:engine}) for both puzzles at every $N$, with a harness that
  re-applies and verifies every solution and an exact adapter for the Santa 2023 puzzles.
\end{itemize}
We also reproduce five known methods: \Mone~\cite{korf1997optimal}, \Mtwo~\cite{kociemba2phase},
\Mthree~\cite{demaine2011algorithms}, \Mfour~\cite{chervov2025cayleypy} and \Msix, the KMCoders pipeline rebuilt from its
write-up~\cite{kmcoders2024santa}. We run three external baselines unchanged (\S\ref{sec:baselines}): \Bone, \Btwo and
\Bthree~\cite{walton2017nxnxn}. Table~\ref{tab:names} lists every name.
